% ==============================================================================
% Section 7: Results
% ==============================================================================

\section{Results}
\label{sec:results}

Results will be reported following the completion of the evaluation study described in \Cref{sec:methodology}. The following subsections outline the planned structure of the results presentation.

\subsection{Vocabulary Acquisition}
\label{subsec:results_vocabulary}

This subsection will report the pre-test and post-test Vocabulary Knowledge Scale (VKS) scores for both the experimental and control groups, including descriptive statistics (means, standard deviations), the mixed-design ANOVA results (F-statistics, p-values, effect sizes), and post-hoc comparisons. Gains will be disaggregated by word difficulty level and domain category.

\subsection{Engagement Metrics}
\label{subsec:results_engagement}

This subsection will present interaction log data characterising usage patterns in the experimental group, including: total number of scanning sessions, mean session duration, words scanned per session, explanation view rates, quiz participation and accuracy rates, gamification engagement (points earned, badges unlocked, streak durations), and usage trajectories over the four-week period. Correlations between gamification engagement and sustained usage will be reported.

\subsection{Learning Outcomes by Prior Knowledge}
\label{subsec:results_outcomes}

This subsection will compare learning gains across prior knowledge subgroups (low, medium, high) within the experimental group, assessing whether the adaptive ITS component produces more equitable outcomes. BKT model accuracy will be evaluated by comparing predicted mastery probabilities with observed quiz performance.

\subsection{Usability}
\label{subsec:results_usability}

This subsection will report System Usability Scale (SUS) scores for the experimental group, including overall scores and item-level analysis. ARCS motivation questionnaire results will be presented for each dimension (Attention, Relevance, Confidence, Satisfaction), with comparison to published norms for mobile learning applications. Key themes from the qualitative interview analysis will be integrated to contextualise the quantitative usability findings.
